\documentclass[11pt,a4paper]{article}
\usepackage[T1]{fontenc}
\usepackage[utf8]{inputenc}
\usepackage{lmodern}
\usepackage{amsmath,amssymb,amsthm,mathtools}
\usepackage[a4paper,margin=27mm,headheight=15pt]{geometry}
\usepackage{microtype,booktabs,tabularx,longtable,array}
\usepackage{xcolor}
\usepackage{enumitem}
\usepackage{fancyhdr}
\usepackage{hyperref}
\usepackage[nameinlink,noabbrev]{cleveref}
\usepackage{listings}
\definecolor{linkblue}{RGB}{25,57,91}
\hypersetup{colorlinks=true,linkcolor=linkblue,citecolor=linkblue,urlcolor=linkblue,
 pdftitle={Spectral Collapse and Alternating RMS Asymptotics for the Rvachev Up-Function},
 pdfauthor={Research manuscript prepared for Vladimir Reshetnikov},
 pdfsubject={Exact analytic transfer spectra, a finite-matrix determinant formula, and certified RMS asymptotics}}
\pagestyle{fancy}
\fancyhf{}
\fancyhead[L]{\small Spectral collapse and Rvachev RMS asymptotics}
\fancyhead[R]{\small Research manuscript}
\fancyfoot[C]{\thepage}
\setlength{\parskip}{0.35em}
\setlength{\parindent}{1.15em}
\setlength{\emergencystretch}{2em}
\setlist[enumerate]{itemsep=0.3em,topsep=0.35em}
\lstset{basicstyle=\ttfamily\small,breaklines=true,columns=fullflexible,frame=single,
 aboveskip=0.7em,belowskip=0.7em}
\newtheorem{theorem}{Theorem}[section]
\newtheorem{lemma}[theorem]{Lemma}
\newtheorem{proposition}[theorem]{Proposition}
\newtheorem{corollary}[theorem]{Corollary}
\theoremstyle{definition}
\newtheorem{definition}[theorem]{Definition}
\newtheorem{example}[theorem]{Example}
\newtheorem{question}[theorem]{Question}
\theoremstyle{remark}
\newtheorem{remark}[theorem]{Remark}
% ed. (2026-09-29): unnumbered environment for editorial notes.
\newtheorem*{ednote}{Editorial note (ProveIt, 2026-09-29)}
\newcommand{\C}{\mathbb C}
\newcommand{\R}{\mathbb R}
\newcommand{\Z}{\mathbb Z}
\newcommand{\N}{\mathbb N}
\newcommand{\dd}{\,\mathrm d}
\newcommand{\ii}{\mathrm i}
\newcommand{\e}{\mathrm e}
\newcommand{\one}{\mathbf 1}
\newcommand{\Lop}{\mathcal L}
\newcommand{\Pop}{\mathcal P}
\newcommand{\Spec}{\operatorname{Spec}}
\newcommand{\tr}{\operatorname{tr}}
\newcommand{\Span}{\operatorname{span}}
\newcommand{\sinc}{\operatorname{sinc}}
\newcommand{\Log}{\operatorname{Log}}
\newcommand{\abs}[1]{\left|#1\right|}
\newcommand{\norm}[1]{\left\|#1\right\|}
\newcommand{\repoSHA}{\texttt{de4ac5ab1438bcbeed58481c4449b593074a1a1f}}
\newcommand{\blobSHA}{\texttt{aeba04f6c5f2d1eccff67645f773ba71803fa2a7}}
\numberwithin{equation}{section}

\title{\textbf{Spectral Collapse and Alternating RMS Asymptotics\\
for the Rvachev Up-Function}\\[0.6em]
\large An exact finite-matrix and boundary-factor theorem\\
for analytic transfer operators}
\author{Research manuscript prepared for Vladimir Reshetnikov}
\date{28 September 2026}

\begin{document}
\maketitle
\begin{abstract}
The Fourier-decay research in the ProveIt repository explicitly leaves open
whether the exhibited eigenvalue $-1/4$ of its quadratic transfer operator is
truly subleading, and whether the observed alternating RMS correction is
nonzero. We resolve these questions on explicitly specified analytic Hardy
spaces. More generally, for every integer dilation $b\ge2$ and trigonometric
polynomial weight $w$, we prove that the Fredholm determinant of the interval
transfer operator is the product of a finite Fourier determinant and the
explicit boundary factor $\prod_{\ell\ge1}(1-zw(0)b^{-\ell})$. When $w(0)=0$,
all nonzero generalized eigenspaces lie in a finite Fourier band.

For $w(x)=\sin^2(\pi x)$ and $b=2$, the complete spectrum is
$\{0,1/2,-1/4\}$, with algebraic multiplicities one and two at the nonzero
values. We derive exact recursions for the spectral projections, a quantitative
remainder of order $\exp(-\tfrac12(\log2)n^2+O(n\log n))$, and a two-term
formula for the squared normalized shell RMS of the Rvachev Fourier product.
An executable certificate using only integer and rational arithmetic proves
that its alternating coefficient is strictly negative. In particular, every
even shell index $n\ge10$ lies below the limiting RMS and every odd index
$n\ge11$ lies above it. We also obtain complete finite-matrix spectral
reductions for all even powers of the sine weight, exact fourth-moment
formulas, and a high-dilation spectral-collapse corollary.

The analytic proofs are conventional mathematical proofs, not new Lean
formalizations. The contribution is a resolution of an identified repository
frontier and explicit extensions; priority for the general determinant
specialization in the wider literature is not asserted.
\end{abstract}
\noindent\textbf{Keywords.} Rvachev up-function; Fabius function; transfer operator;
Fredholm determinant; Thue--Morse product; boundary jets; interval arithmetic.

\clearpage
{\small\setlength{\parskip}{0pt}\tableofcontents}
\clearpage

\section{The question, its scope, and the answer}
\label{sec:intro}

\subsection{A precise repository frontier}
The relevant source is the canonical Fourier-decay synthesis in
ProveIt's Fabius-function project \cite{repo}. Its subsection
\emph{Exact transfer identities without a spectral overclaim} distinguishes
explicit eigenfunctions from a complete spectral theorem. It records that
$-1/4$ is known to be an eigenvalue, but that its subleading status and an
alternating $(-1/2)^n$ relative RMS correction still need a spectral bound on
a specified function space. This is an appropriate target: it has a precise
operator, an observable with concrete asymptotic significance, and a clear
criterion for completion.

The inspected source at commit \repoSHA\ has blob identifier \blobSHA.
The statement of the gap occurs at source lines 1551--1555 of the canonical
TeX file. The full path and immutable link are recorded in
\cref{app:provenance}. The repository was changing during the investigation;
these identifiers, rather than a claim about an indefinitely current branch,
fix the comparison.

There are two logically separate issues. First, an eigenfunction computation
does not exclude other eigenvalues of larger modulus, or additional spectrum
arising from nonperiodic analytic functions. Second, even a genuine
subleading eigenvalue need not occur in the asymptotics of a particular
observable. Its projection can vanish, or integration can annihilate its
eigenfunction. Both issues must be settled to establish a genuine alternating
RMS term.

\subsection{The normalized quadratic operator}
Write
\begin{align}
 (\Lop_2 f)(z)
 &=\frac12\left(\sin^2\frac{\pi z}{2}\,f(z/2)
       +\cos^2\frac{\pi z}{2}\,f((z+1)/2)\right),\label{eq:L2intro}\\
 \Pop&=2\Lop_2.
\end{align}
On real $x\in[0,1]$, $\Pop$ is a positive operator satisfying
$\Pop\one=\one$ and $\norm{\Pop f}_\infty\le\norm f_\infty$.
For $R>1/2$, let $H_R$ be the Hardy Hilbert space of the disk
$D_R=\{z:|z-1/2|<R\}$, defined precisely in \cref{sec:analytic}.

\begin{theorem}[Quadratic spectral resolution]
\label{thm:intro-spectrum}
On every $H_R$, $R>1/2$, the operator $\Lop_2$ is trace class and
\begin{equation}
 \det(I-z\Lop_2)=(1-z/2)(1+z/4)^2.
 \label{eq:L2detintro}
\end{equation}
Its complete spectrum is $\{0,1/2,-1/4\}$. The nonzero eigenvalues are
semisimple, with respective algebraic multiplicities one and two. Their
eigenspaces are
\begin{equation}
 \Span\{\one\},\qquad
 \Span\{u,v\},\quad
 u(x)=\cos(2\pi x)+\tfrac13,\quad v(x)=\sin(2\pi x).
 \label{eq:uv}
\end{equation}
Consequently, continuous linear functionals $\mu,a,b$ on $H_R$ satisfy, for $n\ge1$,
\begin{equation}
 \Pop^n f=\mu(f)\one+(-1/2)^n\bigl(a(f)u+b(f)v\bigr)+Q^nf,
 \label{eq:main-decomp}
\end{equation}
where $Q$ is quasinilpotent and, for every $\varepsilon>0$,
$\norm{Q^n}_{H_R\to H_R}=O_{\varepsilon,R}(\varepsilon^n)$.
\end{theorem}

Here ``quasinilpotent'' means spectral radius zero. It does not mean
nilpotent, finite rank, or identically zero. The statement is on the stated
analytic Hilbert space; it is not a claim about the spectrum on $C[0,1]$ or
$L^p[0,1]$.

\subsection{The RMS conclusion}
Define the entire product and real envelope
\begin{equation}
 \Phi(t)=\prod_{j=0}^{\infty}\frac{\sin(\pi t/2^j)}{\pi t/2^j},\qquad
 E(t)=\exp\!\left(-\frac{(\log t)^2}{2\log2}\right)t^{-\kappa},\quad
 \kappa=\frac{\log(2\pi)}{\log2}.
 \label{eq:PhiE}
\end{equation}
The value of each sinc factor at zero is one. The product is the repository's
Fourier-transform normalization for the Rvachev up-function; its classical
construction is reviewed in \cite{arias}. For $n\ge1$, put
\begin{equation}
 \mathcal R_n^2=\frac{1}{2^{n-1}}
   \int_{2^{n-1}}^{2^n}\abs{\frac{\Phi(t)}{E(t)}}^2\dd t.
 \label{eq:RMSintro}
\end{equation}
The normalization agrees with the source's $\mathcal M_{2,n}$, without a
shift in shell index.

\begin{theorem}[Nonzero alternating RMS correction]
\label{thm:intro-rms}
There are explicitly defined constants $M>0$ and $B<0$ such that
\begin{align}
 \mathcal R_n^2&=\frac M2+\frac B6(-1/2)^n+\epsilon_n,\notag\\
 |\epsilon_n|&\le C_0\exp\!\left(-\tfrac12(\log2)n^2+C_1n\log(n+1)\right).
 \label{eq:RMSmainintro}
\end{align}
Here $C_0,C_1>0$ are constants independent of $n$. With $A=\sqrt{M/2}$,
\begin{equation}
 \mathcal R_n=A+\frac{B}{12A}(-1/2)^n+O(4^{-n}).
 \label{eq:RMSrootintro}
\end{equation}
The supplied integer-arithmetic certificate proves the strict enclosures
\begin{align}
 -0.002907&<B<-0.002898,\label{eq:Bboundsintro}\\
 0.021221672&<M<0.021221677,\label{eq:Mboundsintro}\\
 0.10300890&<A<0.10300894.\label{eq:Aboundsintro}
\end{align}
Furthermore, $\mathcal R_n<A$ for every even $n\ge10$, and
$\mathcal R_n>A$ for every odd $n\ge11$.
\end{theorem}

The constants are $M=\mu(H)$ and $B=a(H)$ for the explicitly given analytic
function $H$ in \cref{eq:H}. The exact coefficient definitions, the analytic
remainder bounds, and the finite certificate are all included below. Thus
\cref{thm:intro-rms} does not infer nonvanishing from a floating-point plot.

\subsection{Relation to established theory and novelty}
Analytic transfer operators, periodic-orbit traces, and dynamical determinants
belong to classical theory; Ruelle's work \cite{ruelle} is a central early
reference. The trace-class determinant facts used here are standard and are
stated explicitly in \cref{sec:analytic}; see \cite{simon,bornemann}.
The relation of the sine product to the Thue--Morse measure and its Riesz
products is also established mathematics \cite{baake}. Related transfer
operators have been studied on continuous and $L^p$ spaces, for example in
\cite{dutkay}; those settings must not be silently identified with the
analytic interval space used here.

The new work relative to the inspected repository consists of the complete
analytic spectral proof, its boundary correction, the exact projection
recursions used for certification, the quantitative remainder, and the
certified nonzero correction for this specific Rvachev observable. The
finite-matrix theorem below is presented with a proof rather than a global
priority claim. No assertion that an internationally recognized longstanding
open problem has been solved is needed or made.

\section{Analytic spaces, compactness, and traces}
\label{sec:analytic}

\subsection{The Hardy disk space}
For $R>1/2$, define
\begin{equation}
 H_R=\left\{f(z)=\sum_{k\ge0}f_k
       \left(\frac{z-1/2}{R}\right)^k:
       \sum_{k\ge0}|f_k|^2<\infty\right\},\qquad
 \norm f_{H_R}^2=\sum_{k\ge0}|f_k|^2.
 \label{eq:Hardy}
\end{equation}
This definition gives a separable Hilbert space with orthonormal basis
$p_k(z)=((z-1/2)/R)^k$. Cauchy--Schwarz yields
\begin{equation}
 |f(z)|\le\frac{\norm f_{H_R}}
 {\sqrt{1-|z-1/2|^2/R^2}}.
 \label{eq:evaluation}
\end{equation}
In particular, restriction to $[0,1]$ is continuous into the uniform norm.
Every function holomorphic on a neighborhood of the closed disk belongs to
$H_R$.

Let $b\ge2$ and let $w$ be an entire, one-periodic trigonometric polynomial.
Set
\begin{equation}
 (\Lop_{b,w}f)(z)=\frac1b\sum_{a=0}^{b-1}
       w\!\left(\frac{z+a}{b}\right)
       f\!\left(\frac{z+a}{b}\right).
 \label{eq:generalL}
\end{equation}
Every inverse branch $\phi_a(z)=(z+a)/b$ satisfies
\begin{equation}
 \abs{\frac{\phi_a(z)-1/2}{R}}
 \le q_R:=\frac1b+\frac{b-1}{2bR}<1
 \quad(z\in\overline D_R).
 \label{eq:branchcontraction}
\end{equation}

\begin{lemma}[Trace-class realization]
\label{lem:traceclass}
The operator $\Lop_{b,w}:H_R\to H_R$ is trace class.
\end{lemma}
\begin{proof}
Consider one branch $Tf=a(z)f(\phi(z))$, with $a$ holomorphic on a
neighborhood of $\overline D_R$ and $\phi$ one of the inverse branches.
If $A_0=\sup_{\overline D_R}|a|$, then
$\norm{Tp_k}_{H_R}\le A_0q_R^k$, since the Hardy norm is bounded by the
supremum on the boundary for these holomorphic functions. The rank-one
expansion
\[
 Tf=\sum_{k\ge0}f_kTp_k
\]
converges in trace norm: the sum of the trace norms is at most
$A_0/(1-q_R)$. A finite sum of these branch operators is trace class.
\end{proof}

\begin{lemma}[Weighted affine-composition trace]
\label{lem:affine-trace}
Suppose $Tf=a(z)f(\phi(z))$ is as above, $\phi$ is affine with slope
$\alpha$, and $z_*$ is its fixed point. Then
\begin{equation}
 \tr T=\frac{a(z_*)}{1-\alpha}.
 \label{eq:affine-trace}
\end{equation}
\end{lemma}
\begin{proof}
Use the normalized coordinate $t=(z-1/2)/R$, so that
$\phi$ becomes $\psi(t)=\alpha t+\gamma$ and
$|\alpha|+|\gamma|<1$. Choose
$|\gamma|/(1-|\alpha|)<\rho<1$. The $k$th diagonal coefficient of $T$ is
\[
 \frac{1}{2\pi\ii}\int_{|t|=\rho}
      a(t)\psi(t)^k t^{-k-1}\dd t,
\]
where $a(t)$ denotes the weight in the normalized coordinate.
Because $|\psi(t)|<|t|$ on this contour, the geometric series is uniformly
convergent there. Summing gives
\[
 \tr T=\frac{1}{2\pi\ii}\int_{|t|=\rho}
       \frac{a(t)}{t-\psi(t)}\dd t
       =\frac{a(t_*)}{1-\alpha}.
\]
The trace-class expansion in \cref{lem:traceclass} justifies summing the
diagonal entries. The residue formula proves the claim.
\end{proof}

\subsection{The standard spectral input}
For clarity, the functional-analysis facts used subsequently are the
following. A compact operator on an infinite-dimensional complex Hilbert
space has only isolated eigenvalues of finite algebraic multiplicity away
from zero. For a trace-class operator $T$, the Fredholm determinant
$D_T(z)=\det(I-zT)$ is entire, its nonzero zeros are the reciprocals of
nonzero eigenvalues with their algebraic multiplicities, and near zero
\begin{equation}
 \log D_T(z)=-\sum_{n\ge1}\frac{z^n}{n}\tr(T^n).
 \label{eq:Fredholm-input}
\end{equation}
Riesz projections separate any finite set of isolated eigenvalues. Finally,
the spectral-radius formula gives
$\lim_n\norm{Q^n}^{1/n}=0$ if $\Spec(Q)=\{0\}$.
These are standard compact-operator and trace-ideal results, not additional
hypotheses about the particular sine product \cite[Ch.~3]{simon}.
The determinant identities are also summarized in
\cite[\S\S2--3, especially (3.1)--(3.4)]{bornemann}.

All subsequent uses check trace class explicitly. In particular, we do not
apply a trace formula to an operator on $C[0,1]$ merely because a finite
collocation matrix has a trace.

\section{The exact finite-matrix and boundary-factor theorem}
\label{sec:general}

Write
\begin{equation}
 w(z)=\sum_{j=-M}^{M}w_j\e^{2\pi\ii jz},\qquad
 c=w(0)=\sum_{j=-M}^{M}w_j,
 \qquad K=\left\lfloor\frac{M}{b-1}\right\rfloor.
 \label{eq:weightdata}
\end{equation}
The coefficients may be complex. Put $w_j=0$ outside the displayed range,
and use the notation $e_k(z)=\e^{2\pi\ii kz}$.

\begin{lemma}[Invariant Fourier band]
\label{lem:band}
The space $V_K=\Span\{e_k:|k|\le K\}$ is invariant under $\Lop_{b,w}$.
Its matrix, with rows and columns indexed by $-K,\ldots,K$, is
\begin{equation}
 A_{b,w}=(w_{br-k})_{-K\le r,k\le K}.
 \label{eq:matrix}
\end{equation}
\end{lemma}
\begin{proof}
Averaging the roots of unity gives the exact identity
\begin{equation}
 \Lop_{b,w}e_k=
 \sum_{\substack{|j|\le M\\b\mid k+j}}w_j e_{(k+j)/b}.
 \label{eq:mode-general}
\end{equation}
If $|k|\le K$, every resulting index has magnitude at most
$\lfloor(K+M)/b\rfloor\le K$. To check the last inequality, write
$M=(b-1)K+s$ with $0\le s<b-1$.
The coefficient of $e_r$ is $w_{br-k}$, proving both assertions.
\end{proof}

\begin{theorem}[Exact determinant factorization]
\label{thm:general-det}
On every $H_R$, $R>1/2$, one has the entire-function identity
\begin{equation}
 \boxed{\displaystyle
 \det(I-z\Lop_{b,w})=
 \det(I-zA_{b,w})
 \prod_{\ell=1}^{\infty}\left(1-\frac{zc}{b^\ell}\right).}
 \label{eq:general-det}
\end{equation}
Thus the nonzero eigenvalues, counted with algebraic multiplicities, are the
nonzero eigenvalues of $A_{b,w}$ together with
$c/b,c/b^2,c/b^3,\ldots$ when $c\ne0$. Coincident values have their
multiplicities added. Zero is also in the spectrum.
\end{theorem}
\begin{proof}
Fix $n\ge1$ and put $B=b^n$ and
\begin{equation}
 W_n(y)=\prod_{j=0}^{n-1}w(b^jy).
 \label{eq:Wn}
\end{equation}
By iterating the inverse branches and using one-periodicity,
\begin{equation}
 (\Lop_{b,w}^nf)(z)=\frac1B\sum_{a=0}^{B-1}
       W_n\!\left(\frac{z+a}{B}\right)
       f\!\left(\frac{z+a}{B}\right).
 \label{eq:iterate}
\end{equation}
The fixed point of the $a$th branch is $a/(B-1)$. Applying
\cref{lem:affine-trace} to all branches yields
\begin{equation}
 \tr(\Lop_{b,w}^n)=\frac{1}{B-1}
    \sum_{a=0}^{B-1}W_n\!\left(\frac{a}{B-1}\right).
 \label{eq:fulltrace}
\end{equation}
Notice that both $0$ and $1$ occur. They are distinct interval points even
though $W_n$ is periodic.

The Fourier degree of $W_n$ is at most
$M(B-1)/(b-1)$. If $W_n=\sum_\ell d_\ell e_\ell$, then the coefficient of
$e_k$ in $\Lop_{b,w}^ne_k$ is $d_{(B-1)k}$. Therefore
\begin{align}
 \tr(A_{b,w}^n)
 &=\sum_{|k|\le K}d_{(B-1)k}\notag\\
 &=\frac{1}{B-1}\sum_{a=0}^{B-2}
      W_n\!\left(\frac{a}{B-1}\right).
 \label{eq:bandtrace}
\end{align}
The roots-of-unity average selects exactly those Fourier coefficients, since
$|(B-1)k|\le M(B-1)/(b-1)$ is equivalent to $|k|\le K$ for integral $k$.
This also covers $B-1=1$.

At the repeated endpoint, $W_n(1)=c^n$. Subtraction gives
\begin{equation}
 \tr(\Lop_{b,w}^n)-\tr(A_{b,w}^n)
 =\frac{c^n}{b^n-1}
 =\sum_{\ell\ge1}\left(\frac{c}{b^\ell}\right)^n.
 \label{eq:trace-difference}
\end{equation}
Use \cref{eq:Fredholm-input} near $z=0$. Absolute convergence there permits
interchanging the two sums, giving
\[
 \log\frac{\det(I-z\Lop_{b,w})}{\det(I-zA_{b,w})}
 =\sum_{\ell\ge1}\log(1-zc/b^\ell).
\]
The infinite product converges locally uniformly on the entire complex plane
because $\sum_\ell |c|b^{-\ell}<\infty$. The two entire functions agree
near zero, and hence everywhere. The spectral conclusions follow from the
standard determinant characterization of algebraic multiplicity. Compactness
on the infinite-dimensional space supplies zero in the spectrum.
\end{proof}

\begin{corollary}[Boundary-zero collapse]
\label{cor:collapse}
If $w(0)=0$, the Fredholm determinant is a polynomial:
\begin{equation}
 \det(I-z\Lop_{b,w})=\det(I-zA_{b,w}).
 \label{eq:collapse}
\end{equation}
There are only finitely many nonzero spectral values. This conclusion is
exact, not a limit of finite-section spectra.
\end{corollary}

\begin{example}[Why the boundary factor cannot be deleted]
\label{ex:constant}
For $w\equiv1$, $K=0$ and $A=[1]$. Nevertheless, on the analytic interval
space the spectrum is
\[
 \{0,1,b^{-1},b^{-2},\ldots\},
\]
not merely $\{0,1\}$. The Bernoulli polynomials give explicit eigenfunctions:
$\Lop_{b,1}B_j=b^{-j}B_j$. Indeed, their generating function
$t\e^{tz}/(\e^t-1)$ transforms under branch averaging into
$(t/b)\e^{(t/b)z}/(\e^{t/b}-1)$. This verifies the eigenidentity coefficient
by coefficient. The endpoint factor in \cref{eq:general-det} is exactly what
accounts for these nonperiodic eigenfunctions.
\end{example}

\section{A direct explanation of spectral collapse}
\label{sec:direct}
The determinant argument proves the complete spectrum and multiplicities.
A second argument explains why no nonperiodic generalized eigenfunction can
survive when the boundary weight vanishes. This is useful both as an
independent check and as a potential formalization route.

\begin{theorem}[Generalized eigenspaces lie in the Fourier band]
\label{thm:generalized-band}
Assume $w(0)=0$ and $\lambda\ne0$. Every $f\in H_R$ satisfying
$(\Lop_{b,w}-\lambda I)^qf=0$ for some $q\ge1$ extends to an entire
one-periodic function and belongs to $V_K$. Consequently its Jordan-chain
structure is exactly the one in the finite matrix $A_{b,w}$.
\end{theorem}
\begin{proof}
We use induction on the length of a generalized eigenvector chain. Let
$g=(\Lop_{b,w}-\lambda I)f$; at the first step $g=0$, and at later steps
$g$ is entire and in $V_K$ by induction.

\emph{Entire extension.} The equation
$f=\lambda^{-1}(\Lop_{b,w}f-g)$ extends $f$ from $D_R$ to
$D_{bR-(b-1)/2}$. Indeed all inverse branches map the latter disk into the
former. Iteration gives radii $1/2+b^n(R-1/2)\to\infty$.
The extensions agree on their overlaps by the original equation and the
identity theorem.

\emph{Periodicity.} For an entire function define
$\Delta f(z)=f(z+1)-f(z)$. Reindex the inverse branches. All intermediate
terms cancel, and periodicity of $w$ gives
\begin{equation}
 \Delta(\Lop_{b,w}f)(z)=\frac1b w(z/b)\Delta f(z/b).
 \label{eq:boundaryintertwine}
\end{equation}
Because $g$ is periodic, $h=\Delta f$ satisfies
\[
 \lambda h(z)=\frac1b w(z/b)h(z/b).
\]
If $h$ were nonzero, let $r$ be its order of vanishing at zero. The left side
has order $r$, while the right side has strictly greater order because
$w(0)=0$. This is impossible. If $w$ is identically zero the conclusion
follows directly as well. Thus $h=0$.

\emph{Finite Fourier support.} An entire periodic function has a Fourier
expansion $f=\sum_{k\in\Z}f_ke_k$, and for every fixed $\tau>0$ there is
$C_\tau$ such that $|f_k|\le C_\tau\e^{-\tau|k|}$. This follows by shifting
the coefficient integral to a horizontal line. Outside $[-K,K]$, the
Fourier coefficients of $g$ vanish, so
\begin{equation}
 \lambda f_r=\sum_{|j|\le M}w_j f_{br-j}\qquad(|r|>K).
 \label{eq:tailrecursion}
\end{equation}
Put $T_N=\sup_{|k|\ge N}|f_k|$ and $C_w=\sum_j|w_j|/|\lambda|$.
For any integer $N>K$,
\[
 T_N\le C_wT_{bN-M}.
\]
The iterated index is
$N_q=b^q(N-M/(b-1))+M/(b-1)$ and tends to infinity geometrically.
Thus $T_N\le C_w^qC_\tau\e^{-\tau N_q}\to0$.
Every Fourier coefficient outside the band vanishes.
\end{proof}

\begin{remark}[Boundary jets explain the geometric ladder]
When $c=w(0)\ne0$, the same difference equation shows that a nonperiodic
eigenfunction with $\Delta f$ vanishing to order $r$ must have eigenvalue
$c/b^{r+1}$. The determinant theorem proves that all these ladder values
occur and gives their multiplicities, including resonances with the finite
Fourier spectrum. It does not assert that resonant Jordan blocks split as
a direct sum without further argument.
\end{remark}

\section{The complete quadratic spectrum and its projections}
\label{sec:quadratic}

\subsection{The three-dimensional calculation}
For $w(x)=\sin^2(\pi x)$, the nonzero Fourier coefficients are
$w_0=1/2$ and $w_{\pm1}=-1/4$, and $w(0)=0$. In the ordered basis
$(e_{-1},e_0,e_1)$,
\begin{equation}
 A_{2,w}=
 \begin{pmatrix}
 -1/4&0&0\\
 -1/4&1/2&-1/4\\
 0&0&-1/4
 \end{pmatrix}.
 \label{eq:A2}
\end{equation}
It follows immediately that
$\det(I-zA_{2,w})=(1-z/2)(1+z/4)^2$. The eigenvectors
$1$, $e_1+1/3$, and $e_{-1}+1/3$ are linearly independent.
Their real and imaginary combinations give $u$ and $v$ in \cref{eq:uv}.
\Cref{thm:general-det,thm:generalized-band} now prove all the spectral
assertions in \cref{thm:intro-spectrum}.

Let $\Pi_0$ and $\Pi_-$ be the Riesz projections of $\Pop$ at $1$ and
$-1/2$. They have the forms
\[
 \Pi_0f=\mu(f)\one,\qquad \Pi_-f=a(f)u+b(f)v.
\]
They commute with $\Pop$ and annihilate one another. For
$Q=\Pop(I-\Pi_0-\Pi_-)$, compact spectral decomposition gives
$\Spec(Q)=\{0\}$ and \cref{eq:main-decomp}. The spectral-radius formula
then gives the asserted faster-than-geometric remainder in the Hardy norm.
This argument does not assume that the operator is normal.

\subsection{Positivity and the invariant probability}
\begin{proposition}
\label{prop:mu}
The functional $\mu$ is integration against a unique stationary probability
measure on $[0,1]$. This measure has full support. In particular,
$\mu(f)>0$ for every nonzero continuous $f\ge0$.
\end{proposition}
\begin{proof}
For real polynomials $f$, \cref{eq:main-decomp} converges uniformly to
$\mu(f)$. Positivity and $\Pop\one=\one$ imply
$|\mu(f)|\le\norm f_\infty$, $\mu(f)\ge0$ when $f\ge0$, and $\mu(1)=1$.
Polynomial density extends $\mu$ to a positive norm-one functional on
$C[0,1]$, represented by a probability measure. Uniform contraction of
$\Pop$ shows, by polynomial approximation, that
$\Pop^nf\to\mu(f)\one$ uniformly for every continuous $f$. Invariance
$\mu(\Pop f)=\mu(f)$ follows; any stationary probability must agree with
this limit and hence with $\mu$.

For full support, let $S$ be the nonempty closed support. If $x\in S$ and a
branch probability at $x$ is positive, the corresponding inverse-branch
image belongs to $S$. This follows by applying stationarity to a nonnegative
continuous function supported near that image. If $S$ initially contained
only endpoints, either endpoint sends positive probability to $1/2$.
Thus $S$ contains an interior point $x$. All inverse branches of every order
then have positive probability, so $S$ contains
$(x+j)/2^n$, $0\le j<2^n$. Their union is dense in $[0,1]$.
\end{proof}

This characterization connects the leading functional to the familiar
Thue--Morse Fourier-coefficient recursion \cite{baake}. Singular continuity
is not needed anywhere in the proof of the RMS theorem.

\subsection{Exact coefficient recursions}
The normalized operator acts on every integral Fourier mode by
\begin{equation}
 \Pop e_{2m}=e_m,\qquad
 \Pop e_{2m+1}=-\tfrac12(e_m+e_{m+1}).
 \label{eq:Pmode}
\end{equation}
Define real sequences $\eta(k)$ and $\alpha(k)$ for $k\ge0$ by
\begin{align}
 &\eta(0)=1,\quad \eta(1)=-\tfrac13,\qquad
 \alpha(0)=0,\quad \alpha(1)=1,\notag\\
 &\eta(2m)=\eta(m),\qquad
   \eta(2m+1)=-\tfrac12\bigl(\eta(m)+\eta(m+1)\bigr),\label{eq:eta}\\
 &\alpha(2m)=-2\alpha(m),\qquad
   \alpha(2m+1)=\alpha(m)+\alpha(m+1).
 \label{eq:alpha}
\end{align}
The recursive clauses are used for indices at least two, avoiding a circular
definition at one. Induction proves
\begin{equation}
 |\eta(k)|\le1,\qquad \alpha(k)\in\Z,
 \qquad |\alpha(k)|\le k.
 \label{eq:sequencebounds}
\end{equation}
For example, $\alpha(1),\ldots,\alpha(8)=1,-2,-1,4,-3,2,3,-8$.

\begin{proposition}[Exact finite-mode reduction]
\label{prop:finite-mode}
For $k\ge1$ and $n\ge\lceil\log_2 k\rceil$,
\begin{equation}
 \Pop^ne_k=\eta(k)+(-1/2)^n\alpha(k)(u+\ii v).
 \label{eq:finite-mode}
\end{equation}
The formula for $-k$ is the complex conjugate. Thus the expansion terminates
exactly after finitely many steps for each trigonometric polynomial.
\end{proposition}
\begin{proof}
At $k=1$, $e_1=-1/3+u+\ii v$ and $u+\ii v$ is an eigenfunction with
value $-1/2$. For larger $k$, use \cref{eq:Pmode} and induction. Every
resulting frequency is at most $\lceil k/2\rceil$ in absolute value.
The constant and eigenfunction coefficients are precisely
\cref{eq:eta,eq:alpha}. The depth bound follows by iterating the halving.
\end{proof}

If $g$ is periodic with sufficiently rapidly decaying Fourier coefficients,
these recursions calculate the projections without collocation:
\begin{align}
 \mu(g)&=\sum_{k\in\Z}\widehat g(k)\eta(|k|),\label{eq:mufourier}\\
 a(g)&=\sum_{k\ne0}\widehat g(k)\alpha(|k|),\label{eq:afourier}\\
 b(g)&=\ii\sum_{k\ne0}\operatorname{sgn}(k)
                \widehat g(k)\alpha(|k|).
 \label{eq:bfourier}
\end{align}
For real $g$, the second expression is
$2\sum_{k\ge1}\alpha(k)\Re\widehat g(k)$. It is essential not to assume
that an arbitrary analytic function on the interval is analytically
periodic. The next section provides the required endpoint regularization
and also justifies these coefficient formulas for the functions used here.

\section{Endpoint smoothing and an explicit supergeometric remainder}
\label{sec:quantitative}

\begin{lemma}[Gain of two matching endpoint jets]
\label{lem:jets}
For every function holomorphic near $[0,1]$,
\begin{equation}
 \Delta(\Pop f)(z)=\sin^2(\pi z/2)\Delta f(z/2)
 \label{eq:jetidentity}
\end{equation}
near zero. Consequently $g_r=\Pop^rf$ has matching derivatives at zero and
one through order $2r-1$.
\end{lemma}
\begin{proof}
The identity is \cref{eq:boundaryintertwine} multiplied by two with $b=2$.
The factor $\sin^2(\pi z/2)$ has a zero of order two. Each iteration
therefore increases the order of the endpoint difference by at least two.
\end{proof}

Let $\Omega_\delta=\{z:\operatorname{dist}(z,[0,1])\le\delta\}$, and
suppose $f$ is holomorphic on a neighborhood of this closed set. Put
$M_f=\sup_{\Omega_\delta}|f|$ and
\begin{equation}
 C_\delta=\frac{\sinh(2\pi\delta)}{2\pi\delta}.
 \label{eq:Cdelta}
\end{equation}
Every inverse branch maps $\Omega_\delta$ into
$\Omega_{\delta/2}$. Moreover,
\[
 \abs{\sin^2(\pi z/2)}+\abs{\cos^2(\pi z/2)}
 =\cosh(\pi\Im z).
\]
Using the product identity
$\prod_{j\ge0}\cosh(t/2^j)=\sinh(2t)/(2t)$ gives, uniformly in $r$,
\begin{equation}
 \sup_{\Omega_\delta}|\Pop^rf|\le C_\delta M_f.
 \label{eq:uniform-complex-bound}
\end{equation}
The product identity follows by repeatedly applying
$\sinh(2t)=2\sinh(t)\cosh(t)$ and taking the small-argument limit.

Cauchy's formula on circles of radius $\delta$, followed by $2r$
integrations by parts, now yields
\begin{equation}
 |\widehat{g_r}(k)|\le
 \frac{C_\delta M_f(2r)!}{(2\pi\delta)^{2r}|k|^{2r}}
 \quad(k\ne0).
 \label{eq:Fourierboundgeneral}
\end{equation}
All boundary terms vanish by \cref{lem:jets}. The periodic extension need
only have these finitely many matching jets; no holomorphic periodic
extension is asserted.

\begin{theorem}[Quantitative remainder]
\label{thm:quantitative}
Suppose $f\in H_R$ and is holomorphic on a neighborhood of
$\Omega_\delta$. For integers $n\ge r\ge2$, put $s=n-r$. Then
\begin{align}
 &\norm{\Pop^nf-\mu(f)\one-(-1/2)^n(a(f)u+b(f)v)}_\infty\notag\\
 &\quad\le
 \frac{C_\delta M_f(2r)!}{(2\pi\delta)^{2r}}
 2^{-s(2r-1)}
 \left(\frac4{2r-1}+\frac8{3(2r-2)}\right).
 \label{eq:explicit-remainder}
\end{align}
In particular, choosing $r=\lfloor n/2\rfloor$ for $n\ge4$ gives
\begin{equation}
 \norm{\Pop^nf-\mu(f)\one-(-1/2)^n(a(f)u+b(f)v)}_\infty
 \le C_{f,\delta}
 \exp\!\left(-\frac{\log2}{2}n^2+O_\delta(n\log n)\right).
 \label{eq:n-square-remainder}
\end{equation}
\end{theorem}
\begin{proof}
For $g=g_r$, \cref{eq:Fourierboundgeneral} implies absolute convergence of
its Fourier series and of all three coefficient sums
\cref{eq:mufourier,eq:afourier,eq:bfourier}. Apply $\Pop^s$ termwise in the
uniform norm. For $|k|\le2^s$ the finite-mode formula is exact. For the
remaining modes, $\norm{\Pop^se_k}_\infty\le1$, $|\eta(k)|\le1$, and
$\norm{u\pm\ii v}_\infty\le4/3$. Their individual residuals have norm
at most $2+(4/3)2^{-s}|k|$.

Put $N=2^s$ and
$D_r=C_\delta M_f(2r)!/(2\pi\delta)^{2r}$. The total residual is bounded by
\begin{align*}
 2D_r\sum_{k>N}k^{-2r}\left(2+\frac{4k}{3N}\right)
 &\le D_rN^{1-2r}
     \left(\frac4{2r-1}+\frac8{3(2r-2)}\right),
\end{align*}
using the integral bounds for the two decreasing power sums.
For fixed $r\ge2$, this residual and its product with $2^s$ tend to zero.
Comparison with the unique leading and subleading terms in the Hardy-space
spectral decomposition identifies the Fourier sums with $\mu(g),a(g),b(g)$.
The commuting projection identities give
$\mu(g)=\mu(f)$ and $(a(g),b(g))=(-1/2)^r(a(f),b(f))$.
This proves \cref{eq:explicit-remainder} without an unjustified
interchange in the Hardy norm.

Finally, $\log((2r)!)=O(r\log r)$ and
$(n-r)(2r-1)=n^2/2+O(n)$ at $r=\lfloor n/2\rfloor$.
Substitution gives \cref{eq:n-square-remainder}.
\end{proof}

\begin{remark}[The Hardy hypothesis can be removed from the scalar estimate]
For a function merely holomorphic near $[0,1]$, define the three coefficients
from the absolutely convergent Fourier sums of $\Pop^2f$, undoing the factor
$(-1/2)^2$ in the last two. The same proof gives convergence and the
coefficient limits. Shifting the iteration shows that the definitions are
independent of the smoothing depth. Hence \cref{eq:explicit-remainder}
holds for all such observables, even if they do not extend to a disk $D_R$.
This extension concerns the uniform asymptotic estimate, not an unstated
spectral theorem on a different Banach space.
\end{remark}

\section{From the Fourier product to the shell observable}
\label{sec:shell}

\subsection{Convergence and a useful complex bound}
On compact subsets of $\C$, the factors of $\Phi$ in \cref{eq:PhiE} are
$1+O(4^{-j})$. The product therefore converges locally uniformly and defines
an entire function. The identity
\[
 \frac{\sin u}{u}=\frac12\int_{-1}^1\e^{\ii tu}\dd t
\]
implies $|\sin u/u|\le\e^{|\Im u|}$. Multiplying the bounds gives
\begin{equation}
 |\Phi(z)|\le\exp(2\pi|\Im z|).
 \label{eq:phibound}
\end{equation}
For $0<z<1$, every real factor is positive and the tail converges to a
positive number. Thus $\Phi(z)>0$ there, while $\Phi(1)=0$.

\subsection{Exact cancellation of the shell scale}
Put $a_0=\log2$. Splitting the infinite product after $n$ factors gives
\begin{equation}
 \Phi(2^{n-1}y)=
 \Phi(y/2)(\pi y)^{-n}2^{-n(n-1)/2}
 \prod_{j=0}^{n-1}\sin(\pi2^jy),\qquad 1\le y\le2.
 \label{eq:shellsplit}
\end{equation}
Define
\begin{equation}
 \boxed{\displaystyle
 H(z)=\frac{\Phi((1+z)/2)^2}{\pi^2}
 \exp\!\left(
 \frac{\Log(1+z)^2}{a_0}+(2\kappa-2)\Log(1+z)\right).}
 \label{eq:H}
\end{equation}
We use the principal logarithm. This function is holomorphic on a
neighborhood of $\overline D_1$, since $\Re(1+z)>0$ there.
In particular $H\in H_1$. It is real and strictly positive on $[0,1)$,
with $H(1)=0$.

Using $2^\kappa=2\pi$, a direct simplification of
\cref{eq:shellsplit} gives, for $x=y-1$,
\begin{equation}
 \abs{\frac{\Phi(2^{n-1}y)}{E(2^{n-1}y)}}^2
 =\frac12 H(x)\prod_{j=0}^{n-1}2\sin^2(\pi2^jx).
 \label{eq:shellnormalized}
\end{equation}
For verification of the normalization, the powers independent of $y$
simplify to
\[
 \pi^{-2n}2^{-n(n-1)+(n-1)^2+2\kappa(n-1)}
 =2^{n-1}\pi^{-2};
\]
the remaining power of $y$ is
$\exp((\log y)^2/a_0+(2\kappa-2)\log y)$.

\begin{proposition}[Exact transfer representation of the RMS]
\label{prop:RMStransfer}
For every $n\ge1$,
\begin{equation}
 \boxed{\displaystyle \mathcal R_n^2=\frac12\int_0^1\Pop^nH(x)\dd x.}
 \label{eq:RMStransfer}
\end{equation}
\end{proposition}
\begin{proof}
Change variables $t=2^{n-1}y$ in \cref{eq:RMSintro} and use
\cref{eq:shellnormalized}. The elementary branch-change identity for
$\Pop$ is
\[
 \int_0^1\Pop^nf(x)\dd x
 =\int_0^1 f(x)\prod_{j=0}^{n-1}2\sin^2(\pi2^jx)\dd x.
\]
It follows either by the substitutions in each inverse branch, or directly
from \cref{eq:iterate} with the normalized weight. Applying it to $H$
proves the result.
\end{proof}

This identity is exact for the infinite sinc product. No finite-product
approximation enters the analytic reduction.

\section{The RMS expansion, all orders, and explicit error bounds}
\label{sec:RMS}
Put
\begin{equation}
 M=\mu(H),\qquad B=a(H),\qquad A=\sqrt{M/2}.
 \label{eq:constants}
\end{equation}
By \cref{prop:mu}, $M>0$. Since
$\int_0^1u=1/3$ and $\int_0^1v=0$, integration of
\cref{eq:main-decomp} immediately yields
\begin{equation}
 \mathcal R_n^2=\frac M2+\frac B6(-1/2)^n+
            \frac12\int_0^1 Q^nH(x)\dd x.
 \label{eq:RMSexactremainder}
\end{equation}
The remainder is faster than every geometric sequence, and
\cref{thm:quantitative} gives the more explicit square-exponential bound
in \cref{eq:RMSmainintro}.

\begin{corollary}[All-orders scalar expansion]
\label{cor:allorders}
For sufficiently large $n$ and every $\varepsilon>0$,
\begin{equation}
 \mathcal R_n=
 A\left(1+\frac{B}{3M}(-1/2)^n\right)^{1/2}
 +O_\varepsilon(\varepsilon^n).
 \label{eq:allordersroot}
\end{equation}
For every fixed $J\ge0$, this gives
\begin{equation}
 \mathcal R_n=
 A\sum_{j=0}^{J}\binom{1/2}{j}
       \left(\frac{B}{3M}\right)^j(-1/2)^{jn}
 +O_J(2^{-(J+1)n}).
 \label{eq:allordersbinomial}
\end{equation}
In particular, if $B\ne0$,
\begin{equation}
 (-2)^n(\mathcal R_n^2-A^2)\longrightarrow B/6,
 \qquad
 (-2)^n(\mathcal R_n-A)\longrightarrow B/(12A).
 \label{eq:scaledlimits}
\end{equation}
\end{corollary}
\begin{proof}
The expression inside the square root tends to one. On a fixed positive
neighborhood of $A^2$, the square-root function is Lipschitz. Thus a
supergeometric error in the squared quantity remains supergeometric after
taking its square root. The binomial expansion gives the finite truncations,
and the first-order term gives the two limits.
\end{proof}

\subsection{A fully explicit bound for the actual observable}
The certificate in the next section uses $\delta=1/4$. We establish here the
analytic constant needed both for certification and for a finite-shell sign
result.

\begin{lemma}[Uniform analytic bound]
\label{lem:Hbound}
On $\Omega_{1/4}$, $|H|<1000$. Moreover, for every $r\ge0$,
$|\Pop^rH|<2000$ on this set. For $g=\Pop^4H$,
\begin{equation}
 |\widehat g(k)|\le \frac{C}{|k|^8}\quad(k\ne0),
 \qquad C=3{,}200{,}000.
 \label{eq:certificate-C}
\end{equation}
\end{lemma}
\begin{proof}
If $z\in\Omega_{1/4}$, then
$\Re(1+z)\ge3/4$, $|1+z|\le9/4$, and $|\Im z|\le1/4$.
The principal logarithm $L=\Log(1+z)$ satisfies $|L|<6/5$:
$|\log|1+z||<1$ and $|\arg(1+z)|<1/3$ already suffice.
Since $\log2>2/3$, $1<\kappa<3$, and \cref{eq:phibound} gives
$|\Phi((1+z)/2)|^2\le\e^{2\pi|\Im z|}$, the modulus of the exponential
factors in $H$ is less than
\[
 \exp\!\left(2+\frac{(6/5)^2}{2/3}+4\cdot\frac65\right)
 =\e^{224/25}<\e^9.
\]
Thus $|H|<\e^9/\pi^2<1000$. The last elementary bound follows, for
example, from $\pi>3$ and $\e<11/4$.

Next $C_{1/4}=\sinh(\pi/2)/(\pi/2)<2$, so
\cref{eq:uniform-complex-bound} gives $|\Pop^rH|<2000$.
Cauchy's estimate yields $\norm{g^{(8)}}_\infty\le8!\,2000\,4^8$.
After eight integrations by parts,
\[
 |\widehat g(k)|\le
 8!\,2000\left(\frac{2}{\pi}\right)^8|k|^{-8}
 <3{,}200{,}000|k|^{-8},
\]
using $\pi>3$. The matching endpoint derivatives required for those
integrations follow from \cref{lem:jets}.
\end{proof}

Taking $r=4$ in the proof of \cref{thm:quantitative}, and using this $C$,
we obtain for every $n\ge4$ the particularly simple bound
\begin{equation}
 \boxed{\displaystyle
 \abs{\mathcal R_n^2-\frac M2-\frac B6(-1/2)^n}
 \le \frac{32C}{63}\,2^{-7(n-4)}.}
 \label{eq:fixed-error}
\end{equation}
Indeed, half of the bracket $4/7+8/(3\cdot6)$ is $32/63$.
Although choosing $r$ with $n$ gives a stronger asymptotic bound,
\cref{eq:fixed-error} is already sufficient to certify the sign from shell
ten onward.

\section{A reproducible certificate for the nonzero coefficient}
\label{sec:certificate}

\subsection{Reduction to finitely many Fourier coefficients}
Let $g=\Pop^4H$. Because the subleading eigenvalue of $\Pop$ is $-1/2$,
\begin{equation}
 B=16\,a(g),\qquad M=\mu(g).
 \label{eq:smoothingconstants}
\end{equation}
For $K\ge1$, truncating \cref{eq:afourier,eq:mufourier} and using
\cref{eq:sequencebounds,eq:certificate-C} gives
\begin{align}
 \abs{B-32\sum_{k=1}^K\alpha(k)\Re\widehat g(k)}
 &\le\frac{16C}{3K^6},\label{eq:tail-a}\\
 \abs{M-\widehat g(0)-2\sum_{k=1}^K\eta(k)\Re\widehat g(k)}
 &\le\frac{2C}{7K^7}.
 \label{eq:tail-mu}
\end{align}
For example, the first tail is bounded by
$32C\sum_{k>K}k^{-7}\le16C/(3K^6)$.

Let $N>2K$, and approximate the coefficients by the exact discrete averages
\begin{equation}
 \widetilde g_N(k)=\frac1N\sum_{j=0}^{N-1}
            g(j/N)\e^{-2\pi\ii kj/N}.
 \label{eq:DFT}
\end{equation}
Absolute Fourier convergence proves the aliasing identity
$\widetilde g_N(k)=\sum_{\ell\in\Z}\widehat g(k+\ell N)$.
For $|k|\le K$, therefore,
\begin{equation}
 |\widetilde g_N(k)-\widehat g(k)|
 \le\frac{4C}{(N-K)^8}.
 \label{eq:alias}
\end{equation}
To see this, use $|k+\ell N|\ge|\ell|(N-K)$ and
$\sum_{\ell\ge1}\ell^{-8}<2$.

Consequently the two finite sums
\begin{align}
 B_{K,N}&=32\sum_{k=1}^K\alpha(k)\Re\widetilde g_N(k),\notag\\
 M_{K,N}&=\widetilde g_N(0)+2\sum_{k=1}^K
                         \eta(k)\Re\widetilde g_N(k)
 \label{eq:finite-estimates}
\end{align}
satisfy the explicit error bounds
\begin{align}
 |B-B_{K,N}|&\le E_B:=
 \frac{16C}{3K^6}+\frac{64CK(K+1)}{(N-K)^8},\label{eq:EB}\\
 |M-M_{K,N}|&\le E_M:=
 \frac{2C}{7K^7}+\frac{4C(2K+1)}{(N-K)^8}.
 \label{eq:EM}
\end{align}
The constants used by the program are
\begin{equation}
 r=4,\quad K=128,\quad N=512,\quad J=28,\quad C=3{,}200{,}000.
 \label{eq:certificateparameters}
\end{equation}
Here $J$ is the number of sinc factors used for a sample evaluation.
The rational error bounds satisfy
$E_B<3.888\cdot10^{-6}$ and $E_M<1.632\cdot10^{-9}$.

\subsection{Enclosing the infinite-product samples}
For $z\in[1/2,1]$ and $J\ge2$, write
$\Phi_J(z)=\prod_{j=0}^{J-1}\sinc(\pi z/2^j)$. The elementary inequalities
$1-u^2/6\le\sin u/u\le1$ for the tail arguments, together with
$\prod(1-a_j)\ge1-\sum a_j$ for $0\le a_j\le1$, give
\begin{equation}
 1-\frac{20}{9}4^{-J}
 \le \frac{\Phi(z)}{\Phi_J(z)}\le1.
 \label{eq:sinctail}
\end{equation}
At an endpoint where the finite product vanishes, the equivalent product
inequality is used rather than division. The estimate follows from
\[
 \sum_{j=J}^\infty\frac{(\pi z/2^j)^2}{6}
 =\frac{2\pi^2z^2}{9}4^{-J}
 \le\frac{20}{9}4^{-J}.
\]
Every evaluation of $H$ uses this tail interval; the program never treats a
finite sinc product as equal to the infinite one.

The inverse-branch formula needed for the samples is
\begin{equation}
 g(x)=\sum_{\ell=0}^{15}
 \left(\prod_{p=0}^3\sin^2(\pi2^py_\ell)\right)H(y_\ell),
 \qquad y_\ell=\frac{x+\ell}{16}.
 \label{eq:g-samples}
\end{equation}
For $x=j/512$, all sample points are dyadic rationals with denominator
$8192$. There are exactly $8192$ evaluations of $H$.
The real parts in \cref{eq:finite-estimates} are calculated by finite cosine
sums, not by an uncertified floating-point FFT.

\subsection{Integer-only elementary-function enclosures}
The primary program \texttt{certify\_integer.py} uses fixed-point intervals
$[l2^{-100},h2^{-100}]$ with integer endpoints. Every addition, multiplication,
and division rounds its rational endpoint result outward by integer floor
or ceiling. Square roots use the integer square-root algorithm. Thus the
certificate requires no floating-point library.

The elementary constants and functions are enclosed as follows. Machin's
identity
\[
 \pi=16\arctan(1/5)-4\arctan(1/239)
\]
is evaluated with 40 alternating-series terms and the first omitted term as
a remainder bound. For $|t|\le4$, sinc is evaluated with 20 terms of
\[
 \sinc t=\sum_{k\ge0}\frac{(-1)^kt^{2k}}{(2k+1)!};
\]
after the first term the absolute terms decrease, so the first omitted term
bounds the remainder. For $y\ge1$, with $q=(y-1)/(y+1)$,
\[
 \log y=2\sum_{k=0}^{L-1}\frac{q^{2k+1}}{2k+1}+R_L,
 \qquad 0\le R_L\le
 \frac{2q^{2L+1}}{(2L+1)(1-q^2)}.
\]
The sample logarithms use $L=36$, and $\log(2\pi)$ uses $L=128$.
Finally, for $0\le x\le4$ the exponential is summed through degree 60;
the positive tail is bounded by its next term divided by $1-x/62$.
All argument-range checks are assertions in the program.

These choices are deliberately generous. The limiting accuracy is set by
the proven Fourier-tail estimate, not by the elementary-function precision.
An independent program, \texttt{certify\_constants.py}, repeats the same
finite reduction with \texttt{mpmath.iv} interval arithmetic. It is a
cross-check, not a dependency of the integer-only certificate.

\subsection{Certificate output and the sign theorem}
The integer-only run produced the following outward decimal enclosures.
The JSON file also records the exact rational endpoints; the displayed
decimals are obtained by integer floor and ceiling.
\begin{center}
\begin{tabular}{@{}lll@{}}
\toprule
Quantity & Lower endpoint & Upper endpoint\\
\midrule
$B$ & $-0.002906388035461419823$ & $-0.002898612708277941829$\\
$M$ & $\phantom{-}0.021221672714945662196$ & $\phantom{-}0.021221675977057311570$\\
$A$ & $\phantom{-}0.103008913970941520154$ & $\phantom{-}0.103008921888002768209$\\
\bottomrule
\end{tabular}
\end{center}
These enclosures imply \cref{eq:Bboundsintro,eq:Mboundsintro,eq:Aboundsintro}.
The wider rational bounds, rather than any noncertified extra digits, are
used in the theorem statements.

It remains to verify the promised finite-shell sign. Since $B<-0.002898$,
the magnitude of its alternating term is greater than
$(0.002898/6)2^{-n}$. At $n=10$, the entirely rational inequality
\begin{equation}
 \frac{32C}{63}\,2^{-42}
 <\frac{0.002898}{6}\,2^{-10}
 \label{eq:signcheck}
\end{equation}
holds. For each increment of $n$, the ratio of the error bound in
\cref{eq:fixed-error} to this lower bound for the alternating term decreases
by $2^{-6}$. The error cannot overturn the sign of the alternating term for any $n\ge10$.
The squared RMS is therefore below $A^2$ at even indices and above it at odd
indices. Taking square roots preserves the comparison. This completes the
proof of \cref{thm:intro-rms}, including nonvanishing, not just an upper bound.

\begin{remark}[What is and is not mechanically checked]
The primary certificate reduces the numerical assertions to explicit
integer computations and elementary series inequalities. Its execution was
checked, and the independent interval implementation agreed. It is not a
Lean or Rocq proof object: the Python implementation and interpreter have
not been verified in a proof assistant. The analytic theorem and all tail
bounds are proved in the manuscript. The supplied code and exact output
make the finite numerical step independently reproducible and auditable.
\end{remark}

\section{Even powers, exact moments, and larger dilations}
\label{sec:evenpowers}

\subsection{Every even sine power has a finite analytic spectrum}
For $m\ge1$, let $\Lop_{2m}=\Lop_{2,w_m}$, where
$w_m(x)=\sin^{2m}(\pi x)$. Its Fourier coefficients are
\begin{equation}
 (w_m)_j=\frac{(-1)^j}{4^m}\binom{2m}{m+j},\qquad |j|\le m.
 \label{eq:evenweight}
\end{equation}
The expression $(-1)^j$ has its usual integral-parity meaning also for
negative $j$.

\begin{corollary}[Complete even-power reduction]
\label{cor:evenpower}
For every $m\ge1$, on every $H_R$ the nonzero spectrum of $\Lop_{2m}$ is
exactly the nonzero spectrum of the $(2m+1)$-dimensional rational matrix
\begin{equation}
 (A_m)_{r,k}=\frac{(-1)^{2r-k}}{4^m}
       \binom{2m}{m+2r-k},\qquad -m\le r,k\le m,
 \label{eq:Am}
\end{equation}
where the binomial coefficient is zero outside its natural range.
The determinant identity and every nonzero algebraic multiplicity are exact.
\end{corollary}
\begin{proof}
Here $M=m$, $K=m$, and $w_m(0)=0$. Apply
\cref{thm:general-det,thm:generalized-band} and substitute
\cref{eq:evenweight}.
\end{proof}

For the integer matrix $B_m=4^mA_m$, exact algebra gives
\begin{align}
 \det(YI-B_1)&=(Y-2)(Y+1)^2,\notag\\
 \det(YI-B_2)&=(Y-1)^2(Y+4)(Y^2-2Y-16),\notag\\
 \det(YI-B_3)&=(Y+1)^2(Y^2+8Y-64)
              (Y^3-10Y^2-96Y+512).
 \label{eq:charpolys}
\end{align}
The accompanying algebra script checks these and the corresponding
polynomials for $m=4,5$. The formulas in \cref{eq:charpolys} are finite
consequences of the displayed matrices; they are not decimal eigenvalue
fits.

\begin{proposition}[A universal outer-row factor]
\label{prop:outerfactor}
For every $m\ge1$, $\det(YI-B_m)$ is divisible by
$(Y-(-1)^m)^2$.
\end{proposition}
\begin{proof}
The row with index $r=m$ has only one nonzero entry, at column $k=m$,
equal to $(-1)^m$. Likewise the row $r=-m$ has only its diagonal entry.
Expanding the characteristic determinant along these two outer rows proves
the factor. This is a statement about algebraic multiplicity. It does not
assert that the corresponding coordinate vectors are eigenvectors, or that
all possible resonances are semisimple.
\end{proof}

\subsection{The quartic spectrum and a cancellation lesson}
For $m=2$, the complete nonzero spectrum is
\begin{equation}
 \frac{1+\sqrt{17}}{16},\quad -\frac14,\quad
 \frac{1-\sqrt{17}}{16},\quad \frac1{16},
 \label{eq:quarticspectrum}
\end{equation}
with multiplicity two at $1/16$ and one at the other values. The leading
spectral radius is $(1+\sqrt{17})/16$, and the next modulus is $1/4$.
Nevertheless that subleading operator eigenvalue is absent from the
following natural scalar moment.

\begin{proposition}[Exact fourth moments]
\label{prop:fourthmoment}
Let
\[
 I_n=\int_0^1\prod_{j=0}^{n-1}\sin^4(\pi2^jx)\dd x,
 \qquad I_0=1.
\]
Then
\begin{equation}
 \sum_{n\ge0}I_nz^n=\frac{1+z/4}{1-z/8-z^2/16},
 \label{eq:fourthGF}
\end{equation}
and, with $s=\sqrt{17}$ and $r_\pm=(1\pm s)/16$,
\begin{equation}
 I_n=\frac{s+5}{2s}r_+^n+\frac{s-5}{2s}r_-^n.
 \label{eq:fourthclosed}
\end{equation}
\end{proposition}
\begin{proof}
The branch-change identity gives $I_n=\int_0^1\Lop_4^n\one$.
Starting from $1$, the iterates stay in
$\Span\{1,\cos(2\pi x)\}$. In that basis the matrix is
\[
 T=\begin{pmatrix}3/8&-1/4\\1/8&-1/4\end{pmatrix},
\]
because $\Lop_4\one=3/8+(1/8)\cos(2\pi x)$ and
$\Lop_4\cos(2\pi x)=-1/4-(1/4)\cos(2\pi x)$.
Integration selects the first coordinate. Thus the generating function is
the $(1,1)$ entry of $(I-zT)^{-1}$, which is \cref{eq:fourthGF}.
Equivalently,
$I_{n+2}=I_{n+1}/8+I_n/16$, with $I_0=1$ and $I_1=3/8$.
Solving that recurrence proves \cref{eq:fourthclosed}.
\end{proof}

Here the $-1/4$ eigenfunction of the full quartic operator is proportional
to $\sin(2\pi x)$, which integration annihilates. The scalar correction in
\cref{eq:fourthclosed} instead has relative ratio
$(1-\sqrt{17})/(1+\sqrt{17})$. This example explains why the separate
nonvanishing certificate for the quadratic Rvachev observable was necessary.
The exact fourth-moment formula is included as a derived illustration,
not as a claim of previously unknown classical moment theory.

\subsection{High-dilation collapse}
\begin{corollary}
\label{cor:highdilation}
Let $m\ge1$ and $b\ge m+2$. On every $H_R$, the operator
$\Lop_{b,w_m}$ has spectrum
\begin{equation}
 \left\{0,\;\frac{1}{4^m}\binom{2m}{m}\right\}.
 \label{eq:highdilationspectrum}
\end{equation}
The nonzero eigenvalue is simple. After division by this value, the powers
converge to a rank-one projection faster than every geometric sequence in
the Hardy operator norm.
\end{corollary}
\begin{proof}
The inequality $b-1>m$ gives $K=0$, and the finite matrix is the scalar
constant Fourier coefficient $4^{-m}\binom{2m}{m}$. The boundary factor
vanishes because $w_m(0)=0$. Apply the determinant theorem and compact
spectral decomposition. The remaining spectral radius is zero.
\end{proof}

This is a genuine simplification from the larger dilation. At $b=m+1$,
the finite band can already have dimension three; the strict threshold in
the statement is not to be replaced by a weak inequality without analysis.

\section{Interpretation, limitations, and a formalization route}
\label{sec:status}

\subsection{What has been resolved}
The operator in the repository is exactly \cref{eq:L2intro}. Its known
pointwise eigen-identities are strengthened to a full spectral theorem on
$H_R$, with algebraic multiplicities and a proof that nonperiodic analytic
functions contribute no additional nonzero spectrum. The general theorem
also identifies precisely when that last assertion would fail: a nonzero
boundary value $w(0)$ creates an infinite geometric spectral ladder.

The source's RMS normalization is reduced to an analytic observable that
belongs to the chosen space. Its $-1/2$ projection survives integration and
has a rigorously enclosed negative coefficient. The rate is therefore a
true leading correction, not simply a possible rate in a norm estimate.
The fixed-shell bound certifies alternating sides of the limit from
$n=10$ onward.

\subsection{What has not been claimed}
The finite spectrum does not describe the operator on every regularity
space. The trace-class disk norm is different from the uniform norm on all
continuous functions, and scalar uniform convergence does not imply a
uniform operator-norm spectral gap on $C[0,1]$. The statement for noninteger
powers $|\sin\pi x|^q$ is not covered by the finite trigonometric-polynomial
argument. The integer-only certificate does not constitute a kernel-checked
Lean proof. Finally, the literature review establishes relevant classical
context and a precise repository gap, not a definitive first-publication
priority for every generalization.

\subsection{Dependency and evidence ledger}
\begin{center}
\small
\begin{tabularx}{\textwidth}{@{}>{\raggedright\arraybackslash}p{0.29\textwidth}>{\raggedright\arraybackslash}X>{\raggedright\arraybackslash}p{0.24\textwidth}@{}}
\toprule
Result & Basis & Status in this package\\
\midrule
Fredholm determinant facts & Classical trace-class spectral theory
\cite{simon,bornemann} & Explicitly stated standard input\\
General factorization & Affine trace lemma and roots-of-unity filtering
& Full conventional proof\\
Quadratic complete spectrum & Exact $3\times3$ matrix and general theorem
& Full conventional proof\\
Projection recursions and remainder & Fourier halving, jet gain, Cauchy bounds
& Full conventional proof\\
RMS transfer identity & Exact sinc-product splitting and branch substitution
& Full conventional proof\\
$B<0$ and numerical enclosures & Proven analytic tails plus integer-only finite sums
& Executed computer-assisted certificate\\
Finite matrices and mode identities & Rational symbolic calculations
& Executed checks; not substitutes for general proofs\\
New Lean/Rocq declarations & None supplied or built
& Not claimed\\
\bottomrule
\end{tabularx}
\end{center}

\subsection{A practical formalization sequence}
The shortest route into the repository begins with the finite algebra, not
with a new spectral library. First formalize \cref{eq:Pmode}, the sequences
$\eta,\alpha$, and the exact mode theorem. The existing
\texttt{RMSTransferEigenfunctions} identities already match the base case.
Next formalize the endpoint-difference identity and its derivative-order
consequence. A verified checker for the finite certificate can be kept
small: fixed-point interval operations, alternating-series remainders,
positive-series tails, and exact rational comparisons suffice.

The analytic bridge then has two possible routes. The spectral route
formalizes trace-class weighted composition, the affine fixed-point trace,
and the determinant theorem. The constructive route formalizes the Fourier
remainder estimate and applies it directly to the shell identity. The latter
can establish the RMS expansion without first building a complete Fredholm
determinant API. Both routes should retain the explicit function-space
hypotheses and the distinction between integer certificates and unverified
floating-point computations.

\section{Further research questions and topics}
\label{sec:questions}
The following questions go beyond what is proved here. They are proposed
research targets, not assertions of open status in the entire literature.

\begin{question}[Sharp speed of the quasinilpotent remainder]
Is the coefficient $\tfrac12\log2$ in the uniform upper bound
\cref{eq:n-square-remainder} optimal for a natural bounded class of analytic
observables? Can the stronger damping visible in repeated endpoint
differences improve it? A useful result would give matching upper and lower
bounds for an explicit nonperiodic observable, rather than only an improved
constant in a general estimate.
\end{question}

\begin{question}[Spectra at finite smoothness]
Determine the complete spectrum or essential spectral radius of $\Lop_2$
on $C^r$, H\"older, Sobolev, and weighted endpoint spaces. Which analytic
eigenvalues remain isolated, and how do unmatched endpoint jets interact
with the noncompact part? The exact boundary intertwining relation provides
a starting point, but the analytic determinant cannot simply be transferred.
\end{question}
% ed. (2026-09-29): the Sobolev part is answered by a later package.
\begin{ednote}
The later repository article
\path{rvachev_up_fourier_decay/Sobolev_Spectral_Disks_Rvachev_Thue_Morse/}
(filed 2026-09-29, unreviewed) answers the Sobolev part. On the periodic
spaces $H^{s,\beta}$ (Sobolev spaces with a logarithmic weight) the
spectrum of $\Lop_2$ is the closed disk of radius
$2^{-s}\sqrt{1+\sqrt{17}}/4$, all of it essential spectrum, together with
$1/2$ and $-1/4$; the value $1/2$ lies outside the disk exactly for
$s>s_0=\tfrac12\log_2\frac{1+\sqrt{17}}4\approx0.1785$, and $-1/4$ exactly
for $s>s_0+1$. In particular $-1/4$ is not an isolated eigenvalue on $L^2$
or $H^1$. On the interval spaces $H^r(0,1)$, $r$ a nonnegative integer, the
same description holds with $s=r$; unmatched endpoint jets add no nonzero
spectrum. H\"older, $C^r$ and weighted endpoint spaces are not treated
there.
\end{ednote}

\begin{question}[Even-power matrix structure]
For the integer matrices $B_m$, classify realness, simplicity, interlacing,
and the ordering of nonzero eigenvalue moduli as $m$ grows. Determine when
the forced outer-row eigenvalue $(-1)^m$ resonates with the interior block
and whether nontrivial Jordan chains occur. The exact polynomials in the
package supply initial data, not an induction proof.
\end{question}

\begin{question}[Large-moment asymptotics]
Use the finite matrices to derive asymptotic expansions for the leading and
subleading eigenvalues as $m\to\infty$. Relate these expansions to the
maximizing periodic orbit of the sine product, and quantify the transition
from finite even moments to the extremal envelope.
\end{question}

\begin{question}[Noninteger powers and the arithmetic mean]
For $q\notin2\N$, establish certified subleading spectra and nonzero
observable corrections for the transfer weight $|\sin\pi x|^q$. The case
$q=1$ is particularly relevant to the repository's arithmetic-mean decay.
A suitable Banach space must accommodate the endpoint regularity and zeros;
finite Fourier truncation alone is no longer exact.
\end{question}

\begin{question}[Boundary perturbations and resonances]
For a trigonometric polynomial family $w_\varepsilon$ with
$w_0(0)=0$ and $w_\varepsilon(0)\ne0$, the determinant theorem explicitly
creates a geometric ladder. Describe the spectral projections and Jordan
structure when one of these ladder values meets a finite-band eigenvalue.
Can one obtain uniform perturbation expansions through a resonance?
\end{question}

\begin{question}[Infinite Fourier weights]
Extend the boundary-factor formula to analytic weights with infinitely many
Fourier coefficients by giving an explicit determinant truncation error.
The finite-band theorem suggests an exact boundary factor and a controlled
infinite periodic component. Such a result could turn heuristic spectral
collocation into an a priori certified method.
\end{question}

\begin{question}[Higher-dimensional boundary corrections]
For expanding integral matrices acting on a cube or parallelepiped,
determine the analogues of the finite Fourier band and the repeated-endpoint
factor. One expects distinct contributions from faces and their
intersections. A correct theorem must distinguish the interval or cube
operator from the corresponding torus operator.
\end{question}

\begin{question}[Arithmetic and regularity of the correction functional]
The integer sequence $\alpha$ has $\alpha(2^j)=(-2)^j$, so its linear growth
bound is sharp. Determine optimal regularity classes on which
$g\mapsto\sum_{k\ne0}\alpha(|k|)\widehat g(k)$ extends continuously,
and identify any useful distributional, refinement, or arithmetic
interpretation of this functional.
\end{question}
% ed. (2026-09-29): answered on the Sobolev scale by an earlier and a later
% package of the repository.
\begin{ednote}
On the Sobolev scale this question is answered in the repository. The
sequence $\alpha$ is $\alpha(k)=B_k(-2)$, the Stern polynomials evaluated
at $-2$, and the functional extends continuously to the Sobolev space $H^s$
of the circle exactly for $s>s_0+1$, with $s_0$ as in the note after
``Spectra at finite smoothness''; both facts are in Part~I of
\path{thue-morse/Thue_Morse_Frontier_Deductions/} (filed 2026-09-05, before
this article, which does not cite it). The later article
\path{rvachev_up_fourier_decay/Sobolev_Spectral_Disks_Rvachev_Thue_Morse/} adds the
logarithmic endpoint: on $H^{s_0+1,\beta}$ the functional is continuous
exactly when $\beta>1/2$. Other regularity classes and further arithmetic
interpretations are not treated.
\end{ednote}

\begin{question}[Verified constants and complete repository integration]
Produce Lean or Rocq proof objects for the finite certificate and the
constructive Fourier remainder theorem. Increase the rigorously enclosed
digits of $A$ and $B$, and certify the earliest shell index at which the
alternating sign persists. The theorem here proves the cutoff ten but does
not claim it is minimal.
\end{question}

\section{Conclusion}
The missing implication in the repository was not an additional numerical
eigenvalue calculation. It was a bridge from a finite Fourier computation to
the full analytic interval operator, followed by a check that the resulting
spectral correction actually occurs in the Rvachev observable. The
boundary-factor theorem supplies the first bridge, and the exact Fourier
projection recursion with integer enclosures supplies the second.

The resulting conclusion is stronger than the originally numerical rate:
there are no other nonzero analytic spectral values in the quadratic case,
the remainder after the two spectral terms has an explicit
square-exponential upper bound, and the normalized RMS is provably on
alternating sides of its limit from shell ten onward. The same mechanism
gives exact finite-dimensional spectral problems for every even sine power
and clarifies when nonperiodic boundary modes must instead be retained.

\appendix
\section{Repository provenance and reproducibility}
\label{app:provenance}

\subsection{Inspected repository source}
Repository: \url{https://github.com/VladimirReshetnikov/ProveIt}.
The comparison uses commit \repoSHA, observed on 28 September 2026 in
Pacific time (29 September in UTC). The canonical article's Git blob is
\blobSHA. Its path is
\begin{quote}\small\ttfamily\raggedright
Analysis/FabiusFunction/docs/\allowbreak
semi-formalized-research-frontiers/drafts/\allowbreak
rvachev\_up\_fourier\_decay/\allowbreak
Rvachev\_Up\_Fourier\_Decay/\allowbreak
Rvachev\_Up\_Fourier\_Decay.tex
\end{quote}
The immutable source link is provided in the bibliography as \cite{repo}.
The sibling README identifies the canonical synthesis and distinguishes it
from a separately supplied rewrite. Source lines 1551--1555 state the
unresolved spectral/RMS claim. The source's explicit eigenfunctions and RMS
power are treated as antecedents, not as new contributions of this article.
No repository file was modified during preparation of this package.

\subsection{Files and commands}
The package contains this source and its compiled PDF, the primary
integer-only certificate, an independent interval-arithmetic cross-check,
an exact algebra checker, their executed outputs, a provenance file, and
build instructions. The minimal numerical verification needs only Python:
\begin{lstlisting}
python certify_integer.py --output integer_certificate.json
\end{lstlisting}
The optional cross-check and finite symbolic checks use the versions listed
in \texttt{requirements.txt}:
\begin{lstlisting}
python certify_constants.py --output constants_certificate.json
python verify_algebra.py --output algebra_certificate.json
\end{lstlisting}
The article has an internal bibliography and no external figure dependencies.
It can be rebuilt by running
\begin{lstlisting}
pdflatex -interaction=nonstopmode -halt-on-error article.tex
pdflatex -interaction=nonstopmode -halt-on-error article.tex
pdflatex -interaction=nonstopmode -halt-on-error article.tex
\end{lstlisting}
All assertions in the checkers are part of the intended execution: do not
run Python with its assertion-disabling optimization flag.

\subsection{Finite checks performed}
The algebra checker verifies the characteristic polynomials of $4^mA_m$
for $1\le m\le5$, the semisimplicity of the quadratic nonzero spectrum,
and the double quartic value $1/16$. It checks 80 finite root-of-unity trace
identities across $2\le b\le5$, $0\le M\le4$, and $1\le n\le4$.
It also verifies the exact finite-mode formula for $1\le k\le512$.
These checks detect indexing and normalization mistakes; the proofs in the
main text, not the finite test ranges, establish the general statements.

\begin{thebibliography}{9}\raggedright
\bibitem{repo}
V.~Reshetnikov, \emph{ProveIt: Rvachev up-function Fourier-decay research},
canonical synthesis, repository source at commit
\texttt{de4ac5ab1438bcbeed58481c4449b593074a1a1f}, inspected
28 September 2026 (Pacific time).
\href{https://github.com/VladimirReshetnikov/ProveIt/blob/de4ac5ab1438bcbeed58481c4449b593074a1a1f/Analysis/FabiusFunction/docs/semi-formalized-research-frontiers/drafts/rvachev_up_fourier_decay/Rvachev_Up_Fourier_Decay/Rvachev_Up_Fourier_Decay.tex}{Immutable canonical TeX source}.

\bibitem{arias}
J.~Arias de Reyna, \emph{An infinitely differentiable function with compact
support: Definition and properties}, arXiv:1702.05442 (2017), English
translation of \emph{Definici\'on y estudio de una funci\'on indefinidamente
diferenciable de soporte compacto}, Rev. Real Acad. Ciencias \textbf{76}
(1982), 21--38. \url{https://arxiv.org/abs/1702.05442}.

\bibitem{ruelle}
D.~Ruelle, \emph{Zeta-functions for expanding maps and Anosov flows},
Invent. Math. \textbf{34} (1976), 231--242.
\href{https://doi.org/10.1007/BF01403069}{doi:10.1007/BF01403069}.

\bibitem{simon}
B.~Simon, \emph{Trace Ideals and Their Applications}, second edition,
Mathematical Surveys and Monographs \textbf{120}, American Mathematical
Society, Providence, 2005.

\bibitem{bornemann}
F.~Bornemann, \emph{On the numerical evaluation of Fredholm determinants},
arXiv:0804.2543, especially \S\S2--3 and formulas (3.1)--(3.4).
\url{https://arxiv.org/abs/0804.2543}.

\bibitem{baake}
M.~Baake, P.~Gohlke, M.~Kesseb\"ohmer, and T.~Schindler,
\emph{Scaling properties of the Thue--Morse measure},
Discrete Contin. Dyn. Syst. \textbf{39} (2019), 4157--4185.
\href{https://doi.org/10.3934/dcds.2019168}{doi:10.3934/dcds.2019168};
\url{https://arxiv.org/abs/1810.06949}.

\bibitem{dutkay}
D.~E.~Dutkay, \emph{The spectrum of the wavelet Galerkin operator},
Integral Equations Operator Theory \textbf{50} (2004), no.~4, 477--487;
arXiv posting (2005), \url{https://arxiv.org/abs/math/0511136}.
\end{thebibliography}
\end{document}
